%! Polynomial Functions and End Behavior — Unit 1, Lesson 1.6 — Group Activity
\documentclass[10pt]{article}
\usepackage{precalculus-article}
\usepackage{precalculus-boxes}
\newcommand{\TallMath}[1]{$\displaystyle #1\rule[-1.4em]{0pt}{3.2em}$}

\begin{document}

\pageheader{Unit 1, Lesson 1.6}{Group Activity}
\namepartnerperiod

\vspace{0.12in}

\begin{scenariobox}[End Behavior Explorer]{sky}
\small
Examine the four polynomial functions below. Use the leading term of each to determine
end behavior and write your answers using limit notation.

\medskip
\renewcommand{\arraystretch}{1.5}
\begin{tabular}{@{} c l c l @{}}
  \textbf{A.} & $f(x) = 2x^5 - 8x^3 + x - 4$ &
  \textbf{B.} & $g(x) = -x^4 + 3x^2 + 10$ \\[2pt]
  \textbf{C.} & $h(x) = -4x^3 + 12x^2 - 9x + 1$ &
  \textbf{D.} & $k(x) = 5x^6 - 100x^4 + 3$ \\
\end{tabular}
\end{scenariobox}

\vspace{0.1in}

\begin{tcolorbox}[colback=white, colframe=black!40,
    title=\textbf{Tier R --- Remediate},
    fonttitle=\bfseries, arc=2mm, left=3mm, right=3mm, top=2mm, bottom=2mm]

Use the \textbf{rules table} from your notes to complete the limit notation for each function.
The leading term has been identified for you.

\smallskip
\renewcommand{\arraystretch}{2.0}
\begin{tabular}{|c|c|c|c|c|}
\hline
\textbf{Function} & \textbf{Leading Term} & \textbf{Degree} & $\lim_{x \to \infty}$ & $\lim_{x \to -\infty}$ \\
\hline
$f(x)$ & $2x^5$   & 5 (odd)  & \blank{2cm} & \blank{2cm} \\
$g(x)$ & $-x^4$   & 4 (even) & \blank{2cm} & \blank{2cm} \\
$h(x)$ & $-4x^3$  & 3 (odd)  & \blank{2cm} & \blank{2cm} \\
$k(x)$ & $5x^6$   & 6 (even) & \blank{2cm} & \blank{2cm} \\
\hline
\end{tabular}
\end{tcolorbox}

\vspace{0.1in}

\begin{tcolorbox}[colback=white, colframe=black!40,
    title=\textbf{Tier A --- Apply},
    fonttitle=\bfseries, arc=2mm, left=3mm, right=3mm, top=2mm, bottom=2mm]

\begin{enumerate}[itemsep=8pt]
  \item For each function, identify the leading term yourself, then write the limit notation.

  \renewcommand{\arraystretch}{2.0}
  \begin{tabular}{|c|c|c|c|}
  \hline
  \textbf{Function} & \textbf{Leading Term} & $\lim_{x \to \infty}$ & $\lim_{x \to -\infty}$ \\
  \hline
  $f(x) = 2x^5 - 8x^3 + x - 4$ & \blank{2cm} & \blank{2cm} & \blank{2cm} \\
  $g(x) = -x^4 + 3x^2 + 10$    & \blank{2cm} & \blank{2cm} & \blank{2cm} \\
  $h(x) = -4x^3 + 12x^2 - 9x + 1$ & \blank{2cm} & \blank{2cm} & \blank{2cm} \\
  $k(x) = 5x^6 - 100x^4 + 3$   & \blank{2cm} & \blank{2cm} & \blank{2cm} \\
  \hline
  \end{tabular}

  \item For function $k(x) = 5x^6 - 100x^4 + 3$, a student notices that $-100x^4$ is a large
        negative number for medium values of $x$ (say $x = 5$). She claims the end behavior
        should be $-\infty$. Explain why she is wrong.

        \vspace{0.2in}
        \writeline\\[4pt]
        \writeline

  \item Without a calculator, predict: for which function(s) does $\lim_{x \to \infty} p(x) = \infty$
        AND $\lim_{x \to -\infty} p(x) = -\infty$? Circle the letter(s).
        \quad \textbf{A} \quad \textbf{B} \quad \textbf{C} \quad \textbf{D}
\end{enumerate}
\end{tcolorbox}

\vspace{0.1in}

\begin{tcolorbox}[colback=white, colframe=black!40,
    title=\textbf{Tier E --- Extend},
    fonttitle=\bfseries, arc=2mm, left=3mm, right=3mm, top=2mm, bottom=2mm]

\begin{enumerate}[itemsep=8pt]
  \item Complete all of Tier A first.

  \item Create your own polynomial function that satisfies ALL of the following:
        \begin{itemize}
          \item Even degree
          \item Negative leading coefficient
          \item Exactly 3 distinct real zeros
        \end{itemize}

        Write your polynomial in factored form: \quad $p(x) = $ \blank{8cm}

        What is the degree of your polynomial? \blank{2cm}

  \item Describe the end behavior of your polynomial using limit notation.

        $\lim_{x \to \infty} p(x) = $ \blank{2.5cm} \qquad $\lim_{x \to -\infty} p(x) = $ \blank{2.5cm}

  \item Explain in one sentence why a polynomial with even degree and negative leading coefficient
        must have the same end behavior on both sides.

        \vspace{0.2in}
        \writeline\\[4pt]
        \writeline
\end{enumerate}
\end{tcolorbox}

\end{document}
